\section{Internet of Things \& Industrial Connectivity}\label{IOT}

The \ac{iiot} refers to the networking of physical devices such as sensors, actuators, and machines with computational and communication infrastructures, enabling monitoring, control, and data-driven optimization~\cite{boyes-the-industrial-internet-of-things, LEE-2015-a-cps-for-indsutry40-based-manufacturing-systems}.
The roots of \ac{iiot} can be traced to earlier machine-to-machine (M2M) communication in manufacturing, but \ac{iiot} expands this concept through standardized communication protocols, cloud and edge connectivity, and scalable integration~\cite{lu-indsutry-40-a-survey-on-technologies}.
Unlike conventional automation systems, \ac{iiot} emphasizes large-scale interoperability, continuous data exchange, and the ability to connect heterogeneous devices across organizational and geographic boundaries.

\subsection{Role in HMLV Manufacturing}

For \ac{hmlv} manufacturing, \ac{iiot} provides capabilities that directly address the challenges of variability, short product cycles, and small batch sizes.
Real-time sensor networks capture operational data that can be used to rapidly detect process deviations or quality issues when switching between product variants~\cite{al-fuqaha-internet-of-things-a-survey-on-enabling}.
Connectivity between machines allows production lines to be reconfigured with minimal downtime, supporting the flexibility required for frequent product changes.
At the same time, traceability across diverse product runs ensures compliance and quality standards are maintained~\cite{zhong-intelligent-manufacturing-in-the-context}.
In practice, \ac{iiot} has been applied in adaptive scheduling, predictive maintenance, and operator decision support systems, all of which are highly relevant to \ac{hmlv} contexts.

\subsection{Key Enabling Technologies and Standards}

The effectiveness of \ac{iiot} depends on several enabling technologies.
At the device level, sensor networks and data pipelines provide the foundation for continuous monitoring and feedback loops~\cite{lu-indsutry-40-a-survey-on-technologies}.
At the communication layer, interoperability standards are essential, The \ac{opc ua} has become a cornerstone for industrial connectivity, offering secure, structured communication across diverse systems.
However, recent studies highlight that performance bottlenecks remain, particularly for latency-sensitive manufacturing processes where rapid response is critical~\cite{ladegourdie-performance-analysis-of-upc-ua}.

Beyond protocol-level interoperability, semantic standardization plays a growing role.
Industry 4.0 standards are increasingly represented in knowledge graphs and linked-data frameworks, which improve the ability of diverse systems to integrate and interpret shared information~\cite{grangel-gonzales-analyzing-a-knowledge-graph}.
Finally, connectivity infrastructures such as Ethernet, industrial fieldbuses, and emerging 5G and edge computing platforms provide the backbone for \ac{iiot} scalability and responsiveness.
Low-latency communication is especially important in \ac{hmlv} manufacturing where production lines may need to adapt in real time to accommodate shifting product configurations~\cite{al-fuqaha-internet-of-things-a-survey-on-enabling}.

\subsection{Challenges and Outlook}

Despite its promise, \ac{iiot} adpotion in \ac{hmlv} manufacturing faces several challenges.
Integrating heterogeneous legacy equipment into \ac{iiot} networks is technically demanding and often requires costly retrofitting.
Data management is another key barrier: vast volumes of sensor and machine data must be stored, processed, and filtered to extract actionable insights~\cite{zhong-intelligent-manufacturing-in-the-context}.
Security and privacy concerns are heightened as connectivity expands across organizational boundaries, creating new attack surfaces~\cite{boyes-the-industrial-internet-of-things}.
Organizationally, manufacturers often face cultural and skills-related barriers in adopting \ac{iiot}, as the shift requires both technical and managerial adaptation.

Looking forward, advances in edge and 5G infrastructures are expected to mitigate some of the latency and scalability issues.
Increasing adoption of semantic interoperability standards will reduce integration complexity, while \ac{ai}-driven analytics layered on top of \ac{iiot} data streams are likely to deliver more powerful decision support.
For \ac{hmlv} manufacturers, these developments could significantly reduce the overhead of frequent product changes and make flexible production systems more cost-effective.
